% Options for packages loaded elsewhere
\PassOptionsToPackage{unicode}{hyperref}
\PassOptionsToPackage{hyphens}{url}
\PassOptionsToPackage{dvipsnames,svgnames,x11names}{xcolor}
%
\documentclass[
  letterpaper,
  DIV=11,
  numbers=noendperiod]{scrartcl}

\usepackage{amsmath,amssymb}
\usepackage{iftex}
\ifPDFTeX
  \usepackage[T1]{fontenc}
  \usepackage[utf8]{inputenc}
  \usepackage{textcomp} % provide euro and other symbols
\else % if luatex or xetex
  \usepackage{unicode-math}
  \defaultfontfeatures{Scale=MatchLowercase}
  \defaultfontfeatures[\rmfamily]{Ligatures=TeX,Scale=1}
\fi
\usepackage{lmodern}
\ifPDFTeX\else  
    % xetex/luatex font selection
\fi
% Use upquote if available, for straight quotes in verbatim environments
\IfFileExists{upquote.sty}{\usepackage{upquote}}{}
\IfFileExists{microtype.sty}{% use microtype if available
  \usepackage[]{microtype}
  \UseMicrotypeSet[protrusion]{basicmath} % disable protrusion for tt fonts
}{}
\makeatletter
\@ifundefined{KOMAClassName}{% if non-KOMA class
  \IfFileExists{parskip.sty}{%
    \usepackage{parskip}
  }{% else
    \setlength{\parindent}{0pt}
    \setlength{\parskip}{6pt plus 2pt minus 1pt}}
}{% if KOMA class
  \KOMAoptions{parskip=half}}
\makeatother
\usepackage{xcolor}
\setlength{\emergencystretch}{3em} % prevent overfull lines
\setcounter{secnumdepth}{-\maxdimen} % remove section numbering
% Make \paragraph and \subparagraph free-standing
\makeatletter
\ifx\paragraph\undefined\else
  \let\oldparagraph\paragraph
  \renewcommand{\paragraph}{
    \@ifstar
      \xxxParagraphStar
      \xxxParagraphNoStar
  }
  \newcommand{\xxxParagraphStar}[1]{\oldparagraph*{#1}\mbox{}}
  \newcommand{\xxxParagraphNoStar}[1]{\oldparagraph{#1}\mbox{}}
\fi
\ifx\subparagraph\undefined\else
  \let\oldsubparagraph\subparagraph
  \renewcommand{\subparagraph}{
    \@ifstar
      \xxxSubParagraphStar
      \xxxSubParagraphNoStar
  }
  \newcommand{\xxxSubParagraphStar}[1]{\oldsubparagraph*{#1}\mbox{}}
  \newcommand{\xxxSubParagraphNoStar}[1]{\oldsubparagraph{#1}\mbox{}}
\fi
\makeatother

\usepackage{color}
\usepackage{fancyvrb}
\newcommand{\VerbBar}{|}
\newcommand{\VERB}{\Verb[commandchars=\\\{\}]}
\DefineVerbatimEnvironment{Highlighting}{Verbatim}{commandchars=\\\{\}}
% Add ',fontsize=\small' for more characters per line
\usepackage{framed}
\definecolor{shadecolor}{RGB}{241,243,245}
\newenvironment{Shaded}{\begin{snugshade}}{\end{snugshade}}
\newcommand{\AlertTok}[1]{\textcolor[rgb]{0.68,0.00,0.00}{#1}}
\newcommand{\AnnotationTok}[1]{\textcolor[rgb]{0.37,0.37,0.37}{#1}}
\newcommand{\AttributeTok}[1]{\textcolor[rgb]{0.40,0.45,0.13}{#1}}
\newcommand{\BaseNTok}[1]{\textcolor[rgb]{0.68,0.00,0.00}{#1}}
\newcommand{\BuiltInTok}[1]{\textcolor[rgb]{0.00,0.23,0.31}{#1}}
\newcommand{\CharTok}[1]{\textcolor[rgb]{0.13,0.47,0.30}{#1}}
\newcommand{\CommentTok}[1]{\textcolor[rgb]{0.37,0.37,0.37}{#1}}
\newcommand{\CommentVarTok}[1]{\textcolor[rgb]{0.37,0.37,0.37}{\textit{#1}}}
\newcommand{\ConstantTok}[1]{\textcolor[rgb]{0.56,0.35,0.01}{#1}}
\newcommand{\ControlFlowTok}[1]{\textcolor[rgb]{0.00,0.23,0.31}{\textbf{#1}}}
\newcommand{\DataTypeTok}[1]{\textcolor[rgb]{0.68,0.00,0.00}{#1}}
\newcommand{\DecValTok}[1]{\textcolor[rgb]{0.68,0.00,0.00}{#1}}
\newcommand{\DocumentationTok}[1]{\textcolor[rgb]{0.37,0.37,0.37}{\textit{#1}}}
\newcommand{\ErrorTok}[1]{\textcolor[rgb]{0.68,0.00,0.00}{#1}}
\newcommand{\ExtensionTok}[1]{\textcolor[rgb]{0.00,0.23,0.31}{#1}}
\newcommand{\FloatTok}[1]{\textcolor[rgb]{0.68,0.00,0.00}{#1}}
\newcommand{\FunctionTok}[1]{\textcolor[rgb]{0.28,0.35,0.67}{#1}}
\newcommand{\ImportTok}[1]{\textcolor[rgb]{0.00,0.46,0.62}{#1}}
\newcommand{\InformationTok}[1]{\textcolor[rgb]{0.37,0.37,0.37}{#1}}
\newcommand{\KeywordTok}[1]{\textcolor[rgb]{0.00,0.23,0.31}{\textbf{#1}}}
\newcommand{\NormalTok}[1]{\textcolor[rgb]{0.00,0.23,0.31}{#1}}
\newcommand{\OperatorTok}[1]{\textcolor[rgb]{0.37,0.37,0.37}{#1}}
\newcommand{\OtherTok}[1]{\textcolor[rgb]{0.00,0.23,0.31}{#1}}
\newcommand{\PreprocessorTok}[1]{\textcolor[rgb]{0.68,0.00,0.00}{#1}}
\newcommand{\RegionMarkerTok}[1]{\textcolor[rgb]{0.00,0.23,0.31}{#1}}
\newcommand{\SpecialCharTok}[1]{\textcolor[rgb]{0.37,0.37,0.37}{#1}}
\newcommand{\SpecialStringTok}[1]{\textcolor[rgb]{0.13,0.47,0.30}{#1}}
\newcommand{\StringTok}[1]{\textcolor[rgb]{0.13,0.47,0.30}{#1}}
\newcommand{\VariableTok}[1]{\textcolor[rgb]{0.07,0.07,0.07}{#1}}
\newcommand{\VerbatimStringTok}[1]{\textcolor[rgb]{0.13,0.47,0.30}{#1}}
\newcommand{\WarningTok}[1]{\textcolor[rgb]{0.37,0.37,0.37}{\textit{#1}}}

\providecommand{\tightlist}{%
  \setlength{\itemsep}{0pt}\setlength{\parskip}{0pt}}\usepackage{longtable,booktabs,array}
\usepackage{calc} % for calculating minipage widths
% Correct order of tables after \paragraph or \subparagraph
\usepackage{etoolbox}
\makeatletter
\patchcmd\longtable{\par}{\if@noskipsec\mbox{}\fi\par}{}{}
\makeatother
% Allow footnotes in longtable head/foot
\IfFileExists{footnotehyper.sty}{\usepackage{footnotehyper}}{\usepackage{footnote}}
\makesavenoteenv{longtable}
\usepackage{graphicx}
\makeatletter
\def\maxwidth{\ifdim\Gin@nat@width>\linewidth\linewidth\else\Gin@nat@width\fi}
\def\maxheight{\ifdim\Gin@nat@height>\textheight\textheight\else\Gin@nat@height\fi}
\makeatother
% Scale images if necessary, so that they will not overflow the page
% margins by default, and it is still possible to overwrite the defaults
% using explicit options in \includegraphics[width, height, ...]{}
\setkeys{Gin}{width=\maxwidth,height=\maxheight,keepaspectratio}
% Set default figure placement to htbp
\makeatletter
\def\fps@figure{htbp}
\makeatother
% definitions for citeproc citations
\NewDocumentCommand\citeproctext{}{}
\NewDocumentCommand\citeproc{mm}{%
  \begingroup\def\citeproctext{#2}\cite{#1}\endgroup}
\makeatletter
 % allow citations to break across lines
 \let\@cite@ofmt\@firstofone
 % avoid brackets around text for \cite:
 \def\@biblabel#1{}
 \def\@cite#1#2{{#1\if@tempswa , #2\fi}}
\makeatother
\newlength{\cslhangindent}
\setlength{\cslhangindent}{1.5em}
\newlength{\csllabelwidth}
\setlength{\csllabelwidth}{3em}
\newenvironment{CSLReferences}[2] % #1 hanging-indent, #2 entry-spacing
 {\begin{list}{}{%
  \setlength{\itemindent}{0pt}
  \setlength{\leftmargin}{0pt}
  \setlength{\parsep}{0pt}
  % turn on hanging indent if param 1 is 1
  \ifodd #1
   \setlength{\leftmargin}{\cslhangindent}
   \setlength{\itemindent}{-1\cslhangindent}
  \fi
  % set entry spacing
  \setlength{\itemsep}{#2\baselineskip}}}
 {\end{list}}
\usepackage{calc}
\newcommand{\CSLBlock}[1]{\hfill\break\parbox[t]{\linewidth}{\strut\ignorespaces#1\strut}}
\newcommand{\CSLLeftMargin}[1]{\parbox[t]{\csllabelwidth}{\strut#1\strut}}
\newcommand{\CSLRightInline}[1]{\parbox[t]{\linewidth - \csllabelwidth}{\strut#1\strut}}
\newcommand{\CSLIndent}[1]{\hspace{\cslhangindent}#1}

\KOMAoption{captions}{tableheading}
\usepackage{amsmath}
\usepackage{amssymb}
\makeatletter
\@ifpackageloaded{tcolorbox}{}{\usepackage[skins,breakable]{tcolorbox}}
\@ifpackageloaded{fontawesome5}{}{\usepackage{fontawesome5}}
\definecolor{quarto-callout-color}{HTML}{909090}
\definecolor{quarto-callout-note-color}{HTML}{0758E5}
\definecolor{quarto-callout-important-color}{HTML}{CC1914}
\definecolor{quarto-callout-warning-color}{HTML}{EB9113}
\definecolor{quarto-callout-tip-color}{HTML}{00A047}
\definecolor{quarto-callout-caution-color}{HTML}{FC5300}
\definecolor{quarto-callout-color-frame}{HTML}{acacac}
\definecolor{quarto-callout-note-color-frame}{HTML}{4582ec}
\definecolor{quarto-callout-important-color-frame}{HTML}{d9534f}
\definecolor{quarto-callout-warning-color-frame}{HTML}{f0ad4e}
\definecolor{quarto-callout-tip-color-frame}{HTML}{02b875}
\definecolor{quarto-callout-caution-color-frame}{HTML}{fd7e14}
\makeatother
\makeatletter
\@ifpackageloaded{caption}{}{\usepackage{caption}}
\AtBeginDocument{%
\ifdefined\contentsname
  \renewcommand*\contentsname{Table of contents}
\else
  \newcommand\contentsname{Table of contents}
\fi
\ifdefined\listfigurename
  \renewcommand*\listfigurename{List of Figures}
\else
  \newcommand\listfigurename{List of Figures}
\fi
\ifdefined\listtablename
  \renewcommand*\listtablename{List of Tables}
\else
  \newcommand\listtablename{List of Tables}
\fi
\ifdefined\figurename
  \renewcommand*\figurename{Figure}
\else
  \newcommand\figurename{Figure}
\fi
\ifdefined\tablename
  \renewcommand*\tablename{Table}
\else
  \newcommand\tablename{Table}
\fi
}
\@ifpackageloaded{float}{}{\usepackage{float}}
\floatstyle{ruled}
\@ifundefined{c@chapter}{\newfloat{codelisting}{h}{lop}}{\newfloat{codelisting}{h}{lop}[chapter]}
\floatname{codelisting}{Listing}
\newcommand*\listoflistings{\listof{codelisting}{List of Listings}}
\makeatother
\makeatletter
\makeatother
\makeatletter
\@ifpackageloaded{caption}{}{\usepackage{caption}}
\@ifpackageloaded{subcaption}{}{\usepackage{subcaption}}
\makeatother

\ifLuaTeX
  \usepackage{selnolig}  % disable illegal ligatures
\fi
\usepackage{bookmark}

\IfFileExists{xurl.sty}{\usepackage{xurl}}{} % add URL line breaks if available
\urlstyle{same} % disable monospaced font for URLs
\hypersetup{
  pdftitle={Multivariate Linear Regression},
  colorlinks=true,
  linkcolor={blue},
  filecolor={Maroon},
  citecolor={Blue},
  urlcolor={Blue},
  pdfcreator={LaTeX via pandoc}}


\title{Multivariate Linear Regression}
\author{}
\date{}

\begin{document}
\maketitle


\subsection{General Principles}\label{general-principles}

To study relationships between multiple continuous independent variables
(e.g., the effect of weight and age on height), we can use a multiple
regression approach. Essentially, we extend
\href{1.\%20Linear\%20Regression\%20for\%20continuous\%20variable.qmd}{Linear
Regression for continuous variable} by adding a regression coefficient
\(\beta_x\) for each continuous variable (e.g., \(\beta_{weight}\) and
\(\beta_{age}\)).

\subsection{Considerations}\label{considerations}

\begin{tcolorbox}[enhanced jigsaw, coltitle=black, colbacktitle=quarto-callout-note-color!10!white, colframe=quarto-callout-note-color-frame, toprule=.15mm, bottomtitle=1mm, left=2mm, opacityback=0, colback=white, title=\textcolor{quarto-callout-note-color}{\faInfo}\hspace{0.5em}{Note}, leftrule=.75mm, toptitle=1mm, titlerule=0mm, arc=.35mm, breakable, opacitybacktitle=0.6, rightrule=.15mm, bottomrule=.15mm]

\begin{itemize}
\item
  We have the same considerations as for the
  \href{1.\%20Linear\%20Regression\%20for\%20continuous\%20variable.qmd}{Regression
  for continuous variable}.
\item
  The model interpretation of the regression coefficients \(\beta_x\) is
  considered for fixed values of the other independent variable(s)'
  regression coefficients---i.e., for a given age, \(\beta_{weight}\)
  represents the expected change in the dependent variable (height) for
  each one-unit increase in weight, holding all other variables (e.g.,
  age) constant.
\end{itemize}

\end{tcolorbox}

\subsection{Example}\label{example}

Below is example code demonstrating Bayesian multiple linear regression
using the Bayesian Inference (BI) package. Data consist of three
continuous variables (\emph{height}, \emph{weight}, \emph{age}), and the
goal is to estimate the effect of \emph{weight} and \emph{age} on
\emph{height}. This example is based on McElreath (2018).

\subsubsection{Python}

\begin{Shaded}
\begin{Highlighting}[]
\ImportTok{from}\NormalTok{ BI }\ImportTok{import}\NormalTok{ bi}

\CommentTok{\# Setup device{-}{-}{-}{-}{-}{-}{-}{-}{-}{-}{-}{-}{-}{-}{-}{-}{-}{-}{-}{-}{-}{-}{-}{-}{-}{-}{-}{-}{-}{-}{-}{-}{-}{-}{-}{-}{-}{-}{-}{-}{-}{-}{-}{-}{-}{-}{-}{-}}
\NormalTok{m }\OperatorTok{=}\NormalTok{ bi(platform}\OperatorTok{=}\StringTok{\textquotesingle{}cpu\textquotesingle{}}\NormalTok{)}

\CommentTok{\# Import Data \& Data Manipulation {-}{-}{-}{-}{-}{-}{-}{-}{-}{-}{-}{-}{-}{-}{-}{-}{-}{-}{-}{-}{-}{-}{-}{-}{-}{-}{-}{-}{-}{-}{-}{-}{-}{-}{-}{-}{-}{-}{-}{-}{-}{-}{-}{-}{-}{-}{-}{-}}
\ImportTok{from}\NormalTok{ importlib.resources }\ImportTok{import}\NormalTok{ files}
\CommentTok{\# Import}
\NormalTok{data\_path }\OperatorTok{=}\NormalTok{ m.load.howell1(only\_path }\OperatorTok{=} \VariableTok{True}\NormalTok{)}
\NormalTok{m.data(data\_path, sep}\OperatorTok{=}\StringTok{\textquotesingle{};\textquotesingle{}}\NormalTok{) }
\NormalTok{m.df }\OperatorTok{=}\NormalTok{ m.df[m.df.age }\OperatorTok{\textgreater{}} \DecValTok{18}\NormalTok{] }\CommentTok{\# Subset data to adults}
\NormalTok{m.scale([}\StringTok{\textquotesingle{}weight\textquotesingle{}}\NormalTok{, }\StringTok{\textquotesingle{}age\textquotesingle{}}\NormalTok{]) }\CommentTok{\# Normalize}

\CommentTok{\# Define model {-}{-}{-}{-}{-}{-}{-}{-}{-}{-}{-}{-}{-}{-}{-}{-}{-}{-}{-}{-}{-}{-}{-}{-}{-}{-}{-}{-}{-}{-}{-}{-}{-}{-}{-}{-}{-}{-}{-}{-}{-}{-}{-}{-}{-}{-}{-}{-}}
\KeywordTok{def}\NormalTok{ model(height, weight, age):}
    \CommentTok{\# Parameter prior distributions}
\NormalTok{    alpha }\OperatorTok{=}\NormalTok{ m.dist.normal(}\DecValTok{0}\NormalTok{, }\FloatTok{0.5}\NormalTok{, name }\OperatorTok{=} \StringTok{\textquotesingle{}alpha\textquotesingle{}}\NormalTok{)    }
\NormalTok{    beta1 }\OperatorTok{=}\NormalTok{ m.dist.normal(}\DecValTok{0}\NormalTok{, }\FloatTok{0.5}\NormalTok{, name }\OperatorTok{=} \StringTok{\textquotesingle{}beta1\textquotesingle{}}\NormalTok{)}
\NormalTok{    beta2 }\OperatorTok{=}\NormalTok{ m.dist.normal(}\DecValTok{0}\NormalTok{, }\FloatTok{0.5}\NormalTok{, name }\OperatorTok{=} \StringTok{\textquotesingle{}beta2\textquotesingle{}}\NormalTok{)}
\NormalTok{    sigma }\OperatorTok{=}\NormalTok{ m.dist.uniform(}\DecValTok{0}\NormalTok{, }\DecValTok{50}\NormalTok{, name }\OperatorTok{=} \StringTok{\textquotesingle{}sigma\textquotesingle{}}\NormalTok{)}
    \CommentTok{\# Likelihood}
\NormalTok{    m.dist.normal(alpha }\OperatorTok{+}\NormalTok{ beta1 }\OperatorTok{*}\NormalTok{ weight }\OperatorTok{+}\NormalTok{ beta2 }\OperatorTok{*}\NormalTok{ age, sigma, obs }\OperatorTok{=}\NormalTok{ height)}

\CommentTok{\# Run MCMC {-}{-}{-}{-}{-}{-}{-}{-}{-}{-}{-}{-}{-}{-}{-}{-}{-}{-}{-}{-}{-}{-}{-}{-}{-}{-}{-}{-}{-}{-}{-}{-}{-}{-}{-}{-}{-}{-}{-}{-}{-}{-}{-}{-}{-}{-}{-}{-}}
\NormalTok{m.fit(model, progress\_bar}\OperatorTok{=}\VariableTok{False}\NormalTok{) }\CommentTok{\# Optimize model parameters through MCMC sampling}

\CommentTok{\# Summary {-}{-}{-}{-}{-}{-}{-}{-}{-}{-}{-}{-}{-}{-}{-}{-}{-}{-}{-}{-}{-}{-}{-}{-}{-}{-}{-}{-}{-}{-}{-}{-}{-}{-}{-}{-}{-}{-}{-}{-}{-}{-}{-}{-}{-}{-}{-}{-}}
\NormalTok{m.summary()}
\end{Highlighting}
\end{Shaded}

\begin{verbatim}
/home/sosa/work/3.12venv/lib/python3.10/site-packages/tqdm/auto.py:21: TqdmWarning:

IProgress not found. Please update jupyter and ipywidgets. See https://ipywidgets.readthedocs.io/en/stable/user_install.html
\end{verbatim}

\begin{verbatim}
BI v 0.0.46 package loaded
jax.local_device_count 32
\end{verbatim}

\begin{longtable}[]{@{}llllllllll@{}}
\toprule\noalign{}
& mean & sd & hdi\_5.5\% & hdi\_94.5\% & mcse\_mean & mcse\_sd &
ess\_bulk & ess\_tail & r\_hat \\
\midrule\noalign{}
\endhead
\bottomrule\noalign{}
\endlastfoot
alpha & 5.17 & 0.49 & 4.39 & 5.94 & 0.01 & 0.01 & 4409.28 & 2652.03 &
1.0 \\
beta1 & 0.19 & 0.50 & -0.52 & 1.09 & 0.01 & 0.01 & 4201.55 & 2920.05 &
1.0 \\
beta2 & -0.03 & 0.49 & -0.80 & 0.75 & 0.01 & 0.01 & 3731.14 & 2847.59 &
1.0 \\
sigma & 49.98 & 0.02 & 49.96 & 50.00 & 0.00 & 0.00 & 3204.17 & 1962.63 &
1.0 \\
\end{longtable}

\subsubsection{R}

\begin{Shaded}
\begin{Highlighting}[]
\FunctionTok{library}\NormalTok{(BayesianInference)}
\NormalTok{m}\OtherTok{=}\FunctionTok{importBI}\NormalTok{(}\AttributeTok{platform=}\StringTok{\textquotesingle{}cpu\textquotesingle{}}\NormalTok{)}

\CommentTok{\# Import Data \& Data Manipulation {-}{-}{-}{-}{-}{-}{-}{-}{-}{-}{-}{-}{-}{-}{-}{-}{-}{-}{-}{-}{-}{-}{-}{-}{-}{-}{-}{-}{-}{-}{-}{-}{-}{-}{-}{-}{-}{-}{-}{-}{-}{-}{-}{-}{-}{-}{-}{-}}
\NormalTok{m}\SpecialCharTok{$}\FunctionTok{data}\NormalTok{(m}\SpecialCharTok{$}\NormalTok{load}\SpecialCharTok{$}\FunctionTok{howell1}\NormalTok{(}\AttributeTok{only\_path =}\NormalTok{ T), }\AttributeTok{sep=}\StringTok{\textquotesingle{};\textquotesingle{}}\NormalTok{)}\CommentTok{\# Import}
\NormalTok{m}\SpecialCharTok{$}\NormalTok{df }\OtherTok{=}\NormalTok{ m}\SpecialCharTok{$}\NormalTok{df[m}\SpecialCharTok{$}\NormalTok{df}\SpecialCharTok{$}\NormalTok{age }\SpecialCharTok{\textgreater{}} \DecValTok{18}\NormalTok{,] }\CommentTok{\# Subset data to adults}
\NormalTok{m}\SpecialCharTok{$}\FunctionTok{scale}\NormalTok{(}\FunctionTok{list}\NormalTok{(}\StringTok{\textquotesingle{}weight\textquotesingle{}}\NormalTok{, }\StringTok{\textquotesingle{}age\textquotesingle{}}\NormalTok{)) }\CommentTok{\# Normalize}
\NormalTok{m}\SpecialCharTok{$}\FunctionTok{data\_to\_model}\NormalTok{(}\FunctionTok{list}\NormalTok{(}\StringTok{\textquotesingle{}weight\textquotesingle{}}\NormalTok{, }\StringTok{\textquotesingle{}height\textquotesingle{}}\NormalTok{, }\StringTok{\textquotesingle{}age\textquotesingle{}}\NormalTok{)) }\CommentTok{\# Send to model (convert to jax array)}

\CommentTok{\# Define model {-}{-}{-}{-}{-}{-}{-}{-}{-}{-}{-}{-}{-}{-}{-}{-}{-}{-}{-}{-}{-}{-}{-}{-}{-}{-}{-}{-}{-}{-}{-}{-}{-}{-}{-}{-}{-}{-}{-}{-}{-}{-}{-}{-}{-}{-}{-}{-}}
\NormalTok{model }\OtherTok{\textless{}{-}} \ControlFlowTok{function}\NormalTok{(height, weight, age)\{}
  \CommentTok{\# Parameter prior distributions}
\NormalTok{  alpha }\OtherTok{=} \FunctionTok{bi.dist.normal}\NormalTok{(}\DecValTok{0}\NormalTok{, }\FloatTok{0.5}\NormalTok{, }\AttributeTok{name =} \StringTok{\textquotesingle{}a\textquotesingle{}}\NormalTok{)}
\NormalTok{  beta1 }\OtherTok{=} \FunctionTok{bi.dist.normal}\NormalTok{(}\DecValTok{0}\NormalTok{, }\FloatTok{0.5}\NormalTok{, }\AttributeTok{name =} \StringTok{\textquotesingle{}b1\textquotesingle{}}\NormalTok{)}
\NormalTok{  beta2 }\OtherTok{=} \FunctionTok{bi.dist.normal}\NormalTok{(}\DecValTok{0}\NormalTok{, }\FloatTok{0.5}\NormalTok{, }\AttributeTok{name =} \StringTok{\textquotesingle{}b2\textquotesingle{}}\NormalTok{)   }
\NormalTok{  sigma }\OtherTok{=} \FunctionTok{bi.dist.uniform}\NormalTok{(}\DecValTok{0}\NormalTok{, }\DecValTok{50}\NormalTok{, }\AttributeTok{name =} \StringTok{\textquotesingle{}s\textquotesingle{}}\NormalTok{)}
  \CommentTok{\# Likelihood}
  \FunctionTok{bi.dist.normal}\NormalTok{(alpha }\SpecialCharTok{+}\NormalTok{ beta1 }\SpecialCharTok{*}\NormalTok{ weight }\SpecialCharTok{+}\NormalTok{ beta2 }\SpecialCharTok{*}\NormalTok{ age, sigma, }\AttributeTok{obs=}\NormalTok{height)}
\NormalTok{\}}

\CommentTok{\# Run MCMC {-}{-}{-}{-}{-}{-}{-}{-}{-}{-}{-}{-}{-}{-}{-}{-}{-}{-}{-}{-}{-}{-}{-}{-}{-}{-}{-}{-}{-}{-}{-}{-}{-}{-}{-}{-}{-}{-}{-}{-}{-}{-}{-}{-}{-}{-}{-}{-}}
\NormalTok{m}\SpecialCharTok{$}\FunctionTok{fit}\NormalTok{(model) }\CommentTok{\# Optimize model parameters through MCMC sampling}

\CommentTok{\# Summary {-}{-}{-}{-}{-}{-}{-}{-}{-}{-}{-}{-}{-}{-}{-}{-}{-}{-}{-}{-}{-}{-}{-}{-}{-}{-}{-}{-}{-}{-}{-}{-}{-}{-}{-}{-}{-}{-}{-}{-}{-}{-}{-}{-}{-}{-}{-}{-}}
\NormalTok{m}\SpecialCharTok{$}\FunctionTok{summary}\NormalTok{() }\CommentTok{\# Get posterior distributions}
\end{Highlighting}
\end{Shaded}

\subsubsection{Julia}

\begin{Shaded}
\begin{Highlighting}[]
\ImportTok{using} \BuiltInTok{BayesianInference}

\CommentTok{\# Setup device{-}{-}{-}{-}{-}{-}{-}{-}{-}{-}{-}{-}{-}{-}{-}{-}{-}{-}{-}{-}{-}{-}{-}{-}{-}{-}{-}{-}{-}{-}{-}{-}{-}{-}{-}{-}{-}{-}{-}{-}{-}{-}{-}{-}{-}{-}{-}{-}}
\NormalTok{m }\OperatorTok{=} \FunctionTok{importBI}\NormalTok{(platform}\OperatorTok{=}\StringTok{"cpu"}\NormalTok{)}

\CommentTok{\# Import Data \& Data Manipulation {-}{-}{-}{-}{-}{-}{-}{-}{-}{-}{-}{-}{-}{-}{-}{-}{-}{-}{-}{-}{-}{-}{-}{-}{-}{-}{-}{-}{-}{-}{-}{-}{-}{-}{-}{-}{-}{-}{-}{-}{-}{-}{-}{-}{-}{-}{-}{-}}
\CommentTok{\# Import}
\NormalTok{data\_path }\OperatorTok{=}\NormalTok{ m.load.}\FunctionTok{howell1}\NormalTok{(only\_path }\OperatorTok{=} \ConstantTok{true}\NormalTok{)}
\NormalTok{m.}\FunctionTok{data}\NormalTok{(data\_path, sep}\OperatorTok{=}\CharTok{\textquotesingle{};\textquotesingle{}}\NormalTok{) }
\NormalTok{m.df }\OperatorTok{=}\NormalTok{ m.df[m.df.age }\OperatorTok{\textgreater{}} \FloatTok{18}\NormalTok{] }\CommentTok{\# Subset data to adults}
\NormalTok{m.}\FunctionTok{scale}\NormalTok{([}\StringTok{"weight"}\NormalTok{, }\StringTok{"age"}\NormalTok{]) }\CommentTok{\# Normalize}

\CommentTok{\# Define model {-}{-}{-}{-}{-}{-}{-}{-}{-}{-}{-}{-}{-}{-}{-}{-}{-}{-}{-}{-}{-}{-}{-}{-}{-}{-}{-}{-}{-}{-}{-}{-}{-}{-}{-}{-}{-}{-}{-}{-}{-}{-}{-}{-}{-}{-}{-}{-}}
\PreprocessorTok{@BI} \KeywordTok{function} \FunctionTok{model}\NormalTok{(height, weight, age)}
    \CommentTok{\# Parameter prior distributions}
\NormalTok{    alpha }\OperatorTok{=}\NormalTok{ m.dist.}\FunctionTok{normal}\NormalTok{(}\FloatTok{0}\NormalTok{, }\FloatTok{0.5}\NormalTok{, name }\OperatorTok{=} \StringTok{"alpha"}\NormalTok{)    }
\NormalTok{    beta1 }\OperatorTok{=}\NormalTok{ m.dist.}\FunctionTok{normal}\NormalTok{(}\FloatTok{0}\NormalTok{, }\FloatTok{0.5}\NormalTok{, name }\OperatorTok{=} \StringTok{"beta1"}\NormalTok{)}
\NormalTok{    beta2 }\OperatorTok{=}\NormalTok{ m.dist.}\FunctionTok{normal}\NormalTok{(}\FloatTok{0}\NormalTok{, }\FloatTok{0.5}\NormalTok{, name }\OperatorTok{=} \StringTok{"beta2"}\NormalTok{)}
\NormalTok{    sigma }\OperatorTok{=}\NormalTok{ m.dist.}\FunctionTok{uniform}\NormalTok{(}\FloatTok{0}\NormalTok{, }\FloatTok{50}\NormalTok{, name }\OperatorTok{=} \StringTok{"sigma"}\NormalTok{)}
    \CommentTok{\# Likelihood}
\NormalTok{    m.dist.}\FunctionTok{normal}\NormalTok{(alpha }\OperatorTok{+}\NormalTok{ beta1 }\OperatorTok{*}\NormalTok{ weight }\OperatorTok{+}\NormalTok{ beta2 }\OperatorTok{*}\NormalTok{ age, sigma, obs }\OperatorTok{=}\NormalTok{ height)}
\KeywordTok{end}

\CommentTok{\# Run mcmc {-}{-}{-}{-}{-}{-}{-}{-}{-}{-}{-}{-}{-}{-}{-}{-}{-}{-}{-}{-}{-}{-}{-}{-}{-}{-}{-}{-}{-}{-}{-}{-}{-}{-}{-}{-}{-}{-}{-}{-}{-}{-}{-}{-}{-}{-}{-}{-}}
\NormalTok{m.}\FunctionTok{fit}\NormalTok{(model)  }\CommentTok{\# Optimize model parameters through MCMC sampling}

\CommentTok{\# Summary {-}{-}{-}{-}{-}{-}{-}{-}{-}{-}{-}{-}{-}{-}{-}{-}{-}{-}{-}{-}{-}{-}{-}{-}{-}{-}{-}{-}{-}{-}{-}{-}{-}{-}{-}{-}{-}{-}{-}{-}{-}{-}{-}{-}{-}{-}{-}{-}}
\NormalTok{m.}\FunctionTok{summary}\NormalTok{() }\CommentTok{\# Get posterior distributions}
\end{Highlighting}
\end{Shaded}

\begin{tcolorbox}[enhanced jigsaw, coltitle=black, colbacktitle=quarto-callout-caution-color!10!white, colframe=quarto-callout-caution-color-frame, toprule=.15mm, bottomtitle=1mm, left=2mm, opacityback=0, colback=white, title=\textcolor{quarto-callout-caution-color}{\faFire}\hspace{0.5em}{Caution}, leftrule=.75mm, toptitle=1mm, titlerule=0mm, arc=.35mm, breakable, opacitybacktitle=0.6, rightrule=.15mm, bottomrule=.15mm]

For R users, if you create the regression coefficient in a single call:

\begin{Shaded}
\begin{Highlighting}[]
\NormalTok{betas }\OperatorTok{=}\NormalTok{ bi.dist.normal(}\DecValTok{0}\NormalTok{, }\FloatTok{0.5}\NormalTok{, name }\OperatorTok{=} \StringTok{\textquotesingle{}regression\_coefficients\textquotesingle{}}\NormalTok{, shape }\OperatorTok{=}\NormalTok{ (}\DecValTok{2}\NormalTok{,))}
\end{Highlighting}
\end{Shaded}

you need to index them starting by 0:

\begin{Shaded}
\begin{Highlighting}[]
\NormalTok{ m$normal(alpha }\OperatorTok{+}\NormalTok{ betas[}\DecValTok{0}\NormalTok{] }\OperatorTok{*}\NormalTok{ weight }\OperatorTok{+}\NormalTok{ betas[}\DecValTok{1}\NormalTok{] }\OperatorTok{*}\NormalTok{ age, sigma, obs}\OperatorTok{=}\NormalTok{height)}
\end{Highlighting}
\end{Shaded}

\end{tcolorbox}

\subsection{Mathematical Details}\label{mathematical-details}

\subsubsection{\texorpdfstring{\emph{Frequentist
formulation}}{Frequentist formulation}}\label{frequentist-formulation}

We model the relationship between the independent variables
\((X_{1i}, X_{2i}, ..., X_{[K,i]})\) and the dependent variable \emph{Y}
using the following equation:

\[
𝑌_i = \alpha +\beta_1  𝑋_{[1,i]} + \beta_2  𝑋_{[2,i]} + ... + \beta_n  𝑋_{[K,i]} + \epsilon_i
\]

Where:

\begin{itemize}
\item
  \(Y_i\) is the dependent variable for observation \emph{i}.
\item
  \(\alpha\) is the intercept term.
\item
  \(X_{[1,i]}\), \(X_{[2,i]}\), \ldots, \(X_{[K,i]}\) are the values of
  the independent variables for observation \emph{i}.
\item
  \(\beta_1\), \(\beta_2\), \ldots, \(\beta_K\) are the regression
  coefficients.
\item
  \(\epsilon_i\) is the error term for observation \emph{i}, and the
  vector of the error terms, \(\epsilon\), are assumed to be independent
  and identically distributed.
\end{itemize}

\subsubsection{\texorpdfstring{\emph{Bayesian
formulation}}{Bayesian formulation}}\label{bayesian-formulation}

In the Bayesian formulation, we define each parameter with
\phantomsection\label{prior}{{priors 🛈}}. We can express the Bayesian
model as follows:

\[
𝑌_i \sim \text{Normal}(\alpha + \sum_{k=1}^K  \beta_k  X_{[K,i]}, σ)
\]

\[
\alpha \sim \text{Normal}(0,1)
\]

\[
\beta_k \sim \text{Normal}(0,1)
\]

\[
σ \sim \text{Uniform}(0, 50)
\]

Where:

\begin{itemize}
\item
  \(Y_i\) is the dependent variable for observation \emph{i}.
\item
  \(\alpha\) is the intercept term, which in this case has a unit-normal
  prior.
\item
  \(\beta_k\) are slope coefficients for the \emph{K} distinct
  independent variables, which also have unit-normal priors.
\item
  \(X_{[1,i]}\), \(X_{[2,i]}\), \ldots, \(X_{[K,i]}\) are the values of
  the independent variables for observation \emph{i}.
\item
  \(\sigma\) is a standard deviation parameter, which here has a Uniform
  prior that constrains it to be positive.
\end{itemize}

\subsection{Reference(s)}\label{references}

\phantomsection\label{refs}
\begin{CSLReferences}{1}{0}
\bibitem[\citeproctext]{ref-mcelreath2018statistical}
McElreath, Richard. 2018. \emph{{Statistical Rethinking: A Bayesian
course with examples in R and Stan}}. Chapman; Hall/CRC.

\end{CSLReferences}




\end{document}
